% https://artofproblemsolving.com/community/c798404h3812886

\documentclass{article}
\usepackage{amsmath, amssymb}

\begin{document}

\section*{Questions English}

\begin{enumerate}

% Answer: D
\item What is the value of $201\times9+20\times19-2\times(0+1+9)$?
\begin{enumerate}
    \item $1409$
    \item $1449$
    \item $2019$
    \item $2169$
    \item $2209$
\end{enumerate}

% Answer: E
\item Which of the following primes is a divisor of $1^{6}+2^{5}+3^{4}+4^{3}+5^{2}+6^{1}$?
\begin{enumerate}
    \item $2$
    \item $3$
    \item $5$
    \item $7$
    \item $11$
\end{enumerate}

% Answer: A
\item Let $a$, $b$, $c$ and $d$ be distinct integers between $1$ and $50$ inclusive. How many possible values of $a+b+c+d$ exist?
\begin{enumerate}
    \item $185$
    \item $191$
    \item $196$
    \item $197$
    \item $200$
\end{enumerate}

% Answer: D
\item Let $N$ be the unique positive integer such that the $N$\% of $880$ is a perfect square smaller than $880$. What is the sum of the digits of $N$?
\begin{enumerate}
    \item $7$
    \item $8$
    \item $9$
    \item $10$
    \item $11$
\end{enumerate}

% Answer: B
\item The ratio of the area of an octagon to its perimeter is $5/4$, when calculations are done in inches. When calculations are done in feet, what is the ratio of the area of the same octagon to its perimeter?
\begin{enumerate}
    \item $\frac{5}{576}$
    \item $\frac{5}{48}$
    \item $\frac{5}{4}$
    \item $15$
    \item $180$
\end{enumerate}

% Answer: B
\item Two different numbers are randomly selected from the set $S=\{1,2,3,4,5,6,7,8,9,10,11,...,2019\}$. The probability that their sum is $2020$ would be greater if the number $n$ had first been removed from set $S$. What is the sum of the digits of $n$?
\begin{enumerate}
    \item $1$
    \item $2$
    \item $3$
    \item $10$
    \item $12$
\end{enumerate}

% Answer: D
\item Ten gangsters are standing on a flat surface with the distances between them all distinct. At twelve o'clock, when the church bells start chiming, each of them fatally shoots the one among the other $9$ gangsters who is the farthest. At most how many gangsters will remain alive?
\begin{enumerate}
    \item $1$
    \item $2$
    \item $7$
    \item $8$
    \item $9$
\end{enumerate}

% Answer: D
\item Rectangle $ABCD$ is drawn with $AB=8$ and $BC=6$. The incircle of $\triangle ACD$ has center $I$. What is $BI^{2}$?
\begin{enumerate}
    \item $32$
    \item $36$
    \item $48$
    \item $52$
    \item $64$
\end{enumerate}

% Answer: D
\item We define a real number $x$ to be a semi-integer if $\sqrt{2}\cdot x$ is an integer. How many real numbers $0\le x\le100$ are semi-integers?
\begin{enumerate}
    \item $71$
    \item $140$
    \item $141$
    \item $142$
    \item $211$
\end{enumerate}

% Answer: D
\item What is the sum of all positive integers $n$ such that $3^{20}<n^{n}<8^{20}$?
\begin{enumerate}
    \item $60$
    \item $65$
    \item $66$
    \item $75$
    \item $84$
\end{enumerate}

% Answer: C
\item For real numbers $m$ and $n$, we define the operator $\circ$ as $m\circ n = m^{-1}-(mn)^{-1}+n^{-1}$. What is the value of $1\circ(2\circ(\cdot\cdot\cdot(2018\circ2019)))$?
\begin{enumerate}
    \item $\frac{1}{2019}$
    \item $\frac{2018}{2019}$
    \item $1$
    \item $\frac{2019}{2018}$
    \item $2019$
\end{enumerate}

% Answer: B
\item Regular hexagon $ABCDEF$ has side length $1$. Points $P$ and $Q$ are placed in the interior of the hexagon such that $ABPQ$ is a square. If $M$ and $N$ are the midpoints of sides $AB$ and $DE$, respectively, what is the area of the shaded region?
\begin{enumerate}
    \item $\frac{2}{3}$
    \item $\frac{\sqrt{3}}{2}$
    \item $\frac{8}{9}$
    \item $1$
    \item $\frac{2\sqrt{3}}{3}$
\end{enumerate}

% Answer: C
\item How many ordered triples of primes $(p,q,r)$ satisfy $p^{q}+q^{p}=r$?
\begin{enumerate}
    \item $0$
    \item $1$
    \item $2$
    \item $3$
    \item $6$
\end{enumerate}

% Answer: C
\item Suppose $\min(N)$ is the smallest positive integer such that $\frac{0.25^{1000}}{0.1^{N}}>1$. In which of the following intervals does $\min(N)$ lie between?
\begin{enumerate}
    \item $[250, 350]$
    \item $[400, 500]$
    \item $[550, 650]$
    \item $[700, 800]$
    \item $[850, 950]$
\end{enumerate}

% Answer: D
\item Consider a polynomial $p(x)$ of degree $1$ such that for a real number $a$, $p(a)=2$, $p(p(a))=17$ and $p(p(p(a)))=167$. What is the value of $a$?
\begin{enumerate}
    \item $-1$
    \item $-\frac{1}{2}$
    \item $0$
    \item $\frac{1}{2}$
    \item $1$
\end{enumerate}

% Answer: D
\item Geoff the frog is standing at the origin in the coordinate plane. For each move, Geoff can only move $1$ unit to the right or $1$ unit upwards; also, every up move must be immediately followed by a right move (except for the last move). What is the number of distinct sequences of moves that end at the point $(9,5)$?
\begin{enumerate}
    \item $56$
    \item $126$
    \item $210$
    \item $252$
    \item $346$
\end{enumerate}

% Answer: C
\item In triangle $\triangle ABC$, let $D$ be the midpoint of $BC$ and $E$ be a point on $AD$ so that $\angle BEC+\angle BAC=180^{\circ}$. If $BC=\sqrt{7}$, the value of $DE \cdot DA$ can be written as $\frac{m}{n}$ for some relatively prime integers $m$, $n$. What is $m+n$?
\begin{enumerate}
    \item $7$
    \item $9$
    \item $11$
    \item $51$
    \item $53$
\end{enumerate}

% Answer: D
\item A random number generator generates a random number each second. On the $n$th second, the generator randomly outputs an integer from the set $\{1,2,..., n\}$. What is the expected number of seconds it will take before the sum of the outputted numbers is at least $280$?
\begin{enumerate}
    \item $23$
    \item $24$
    \item $30$
    \item $32$
    \item $35$
\end{enumerate}

% Answer: C
\item Let $a$, $b$, $c$, $d$ be $4$ not necessarily distinct integers chosen from the set $\{0,1,..., 2019\}$. What is the probability that $(a+b)(c+d)+(a+c)(b+d)+(a+d)(b+c)$ is divisible by $4$?
\begin{enumerate}
    \item $\frac{1}{4}$
    \item $\frac{3}{16}$
    \item $\frac{3}{8}$
    \item $\frac{7}{16}$
    \item $\frac{1}{2}$
\end{enumerate}

% Answer: C
\item Let $\triangle ABC$ be an equilateral triangle with its $3$ vertices on the graph $y=|2x|$. If the midpoint of $AC$ is $(3,6)$, then what is the area of $\triangle ABC$?
\begin{enumerate}
    \item $20\sqrt{3}$
    \item $25\sqrt{3}$
    \item $\frac{80\sqrt{3}}{3}$
    \item $\frac{85\sqrt{3}}{3}$
    \item $45\sqrt{3}$
\end{enumerate}

% Answer: C
\item James must travel from point $A$ to point $B$ along the grid shown below. He can only move along the arrows and he must take exactly $8$ steps. How many possible paths are there?
\begin{enumerate}
    \item $70$
    \item $72$
    \item $140$
    \item $360$
    \item $560$
\end{enumerate}

% Answer: D
\item Triangle $\triangle ABC$ has $AB=11$, $BC=13$ and $AC=20$. A circle is drawn with diameter $AC$. Line $AB$ intersects the circle at $D\ne A$ and line $BC$ intersects the circle at $E\ne B$. What is the length of $DE$?
\begin{enumerate}
    \item $\frac{55}{13}$
    \item $5$
    \item $\frac{66}{13}$
    \item $\frac{100}{13}$
    \item $\frac{120}{13}$
\end{enumerate}

% Answer: B
\item For positive integers $a$ and $b$, let $*$ be an operator with the following properties: $a*b=b*a$, $(a+1)*b=a*b+a+b$, $1*1=2$. Let $c$ be the unique positive integer such that there exists exactly $25$ ordered pairs of positive integers $(a, b)$ satisfying $a*b=c$. What is the value of $c$?
\begin{enumerate}
    \item $301$
    \item $326$
    \item $351$
    \item $376$
    \item $401$
\end{enumerate}

% Answer: D
\item Let $\Gamma_{1}$, $\Gamma_{2}$, $\Gamma_{3}$ be $3$ circles with radii $2$, $3$, $6$ respectively, such that $\Gamma_{1}$ and $\Gamma_{2}$ are externally tangent at $A$ and $\Gamma_{3}$ is internally tangent to $\Gamma_{1}$ and $\Gamma_{2}$ at $B$, $C$ respectively. What is the radius of the circumcircle of $\triangle ABC$?
\begin{enumerate}
    \item $\frac{6}{\sqrt{11}}$
    \item $5$
    \item $\sqrt{2}+\sqrt{3}+\sqrt{6}$
    \item $6$
    \item $4\sqrt{3}$
\end{enumerate}

% Answer: A
\item Let $f(x)=\lfloor x\rfloor(x+\{x\})$ be a function defined over all positive reals $x$, where $\lfloor x\rfloor$ denotes the greatest integer less than or equal to $x$ and $\{x\}=x-\lfloor x\rfloor$ denotes the fractional part of $x$. It is given that there exist unique real numbers $\alpha$ and $\beta$ such that $f(\alpha)=343$ and $f(\beta)=334$. What is $\alpha-\beta$?
\begin{enumerate}
    \item $\frac{1}{4}$
    \item $\frac{1}{3}$
    \item $\frac{3}{8}$
    \item $\frac{5}{12}$
    \item $\frac{4}{9}$
\end{enumerate}

\end{enumerate}

\section*{Solutions English}

\begin{enumerate}
\item \textbf{Answer (D):} We compute: $201\times9+20\times19-2\times(0+1+9)=1809+380-20=2169$.

\item \textbf{Answer (E):} We compute: $1^{6}+2^{5}+3^{4}+4^{3}+5^{2}+6^{1}=1+32+81+64+25+6=209$. Since $209=11\cdot19$, $11$ is a prime factor of the expression.

\item \textbf{Answer (A):} The smallest possible value of $a+b+c+d$ is $1+2+3+4=10$, while the largest possible value is $47+48+49+50=194$. Consider the expression $1+2+3+4$. We can keep increasing the $4$ by $1$ until we reach $1+2+3+50$. Next, we can keep increasing the $3$ by $1$ until we reach $1+2+49+50$. Then, we can keep increasing the $2$ by $1$ until we reach $1+48+49+50$. Finally, we can keep increasing the $1$ by $1$ until we reach $47+48+49+50$. Therefore, all integer values between $10$ and $194$ inclusive are obtainable, so there are $194-10+1=185$ possible values.

\item \textbf{Answer (D):} The number that is $N$\% of $880$ is $\frac{N}{100}\cdot880\Rightarrow\frac{44N}{5}$. Since this expression must be an integer, we let $N=5k$ for some positive integer $k$. Then, our expression becomes $44k=2^{2}\cdot11\cdot k$. For this to be a perfect square, the exponent of each prime in the prime factorization must be even. Therefore, $k=11$ and $N=55$. We can check that $55$\% of $880$ is $484=22^{2}$. The requested sum is $5+5=10$. To show that $N$ is unique, we note that for the resulting number to be smaller than $880$, we must have $N<100$ or $k<20$. Since the resulting number is a perfect square if $2^{2}\cdot11\cdot k$ is a perfect square, $k$ must be divisible by $11$. But $11$ is the only positive integer less than $20$ that is divisible by $11$, so $N$ cannot equal another value less than $100$. Thus, $N$ is unique.

\item \textbf{Answer (B):} The ratio of the octagon's area to perimeter in inches is given by $\frac{5~in.^{2}}{4~in.}=\frac{5~in.}{4}$. We note that $1~ft=12~in.$. Converting the unit of measurement to feet: $\frac{5~in.}{4}\cdot\frac{1~ft}{12~in.}=\frac{5~ft}{48}$. Therefore, the requested ratio is $\frac{5}{48}$.

\item \textbf{Answer (B):} First, select one of the $2019$ integers at random and then select one of the remaining $2018$ integers at random. Note that for every positive integer $k$ from $1$ to $2019$ except $1010$, the integer $2020-k$ can be selected after selecting $k$. Therefore, if we select $k\ne1010$ as the first number, there is a $\frac{1}{2018}$ chance that we will select $2020-k$ as the second number. However, if we select $k=1010$ as the first number, there is no chance that we will select $1010$ again, since the two integers must be distinct. Therefore, $n=1010$ and the requested sum is $1+0+1+0=2$.

\item \textbf{Answer (D):} Objectively, to maximize the number of gangsters alive, one gangster should be shot by as many gangsters as possible to prevent other gangsters from dying. Clearly, all $10$ of them cannot remain alive. If $9$ gangsters remain alive, then one gangster must be shot by every gangster, including himself. However, a gangster clearly cannot shoot himself because he is not the gangster furthest away from himself. However, $8$ gangsters can remain alive. This can be done by placing $9$ of the gangsters close together and placing the $10$th gangster arbitrarily far from the other gangsters. Then, the $10$th gangster will shoot one of the first $9$ gangsters, while the first $9$ gangsters will all shoot the $10$th. Therefore, at most $8$ gangsters can remain alive.

\item \textbf{Answer (D):} Clearly, $AC=10$ by Pythagorean Theorem on $\triangle ACD$. Consider the incircle of $\triangle ACD$. Its radius is equal to the area of $\triangle ACD$ divided by its semiperimeter. The area of $\triangle ACD$ is $\frac{1}{2}\cdot6\cdot8=24$. Since the semiperimeter is $12$, the radius of the incircle is $2$. Therefore, the distance from $I$ to $DC$ is $2$ and the distance from $I$ to $AB$ is $6-2=4$. The distance from $I$ to $AD$ is $2$, so the distance from $I$ to $BC$ is $8-2=6$. Therefore, by Pythagorean Theorem, $BI^{2}=4^{2}+6^{2}=52$.

\item \textbf{Answer (D):} For $\sqrt{2}\cdot x$ to be an integer, $x$ must be in the form $\frac{n}{\sqrt{2}}$ for some integer $n$. Then, $0\le x\le100\Rightarrow0\le\frac{n}{\sqrt{2}}\le100\Rightarrow0\le n\le100\sqrt{2}$. The value of $\sqrt{2}$ is often approximated as $1.414$, so $141<100\sqrt{2}<142\Rightarrow0\le n\le141$. There are $142$ possible values of $n$, which gives us $142$ possible values of $x$.

\item \textbf{Answer (D):} Note that $3^{20}=9^{10}$ and $8^{20}=16^{15}$. Our inequality becomes $9^{10}<n^{n}<16^{15}$. Clearly, the values of $n$ that satisfy this inequality are $10, 11, 12, 13, 14$, and $15$, which have a sum of $75$.

\item \textbf{Answer (C):} Note that for all real numbers $k$, $1\circ k=1-\frac{1}{k}+\frac{1}{k}=1$. Therefore, $1\circ(2\circ(\cdot\cdot\cdot(2018\circ2019)))=1$.

\item \textbf{Answer (B):} The shaded area is given by the area of kite $NCMF$ minus the area of kite $NPMQ$. Since the area of a kite is half of the product of its diagonals, our desired area is $\frac{1}{2}(NM\cdot CF-NM\cdot PQ)$. Clearly, $NM=\sqrt{3}$ and $CF=2$. In addition, $PQ=1$ because $ABPQ$ is a square. Therefore, the desired area is $\frac{1}{2}(\sqrt{3}\cdot2-\sqrt{3}\cdot1)=\frac{\sqrt{3}}{2}$.

\item \textbf{Answer (C):} Note that $r=p^{q}+q^{p}\ge2^{2}+2^{2}=8$. Therefore, $r\ne2$, so $r$ must be an odd prime. Since the exponentiation of an integer preserves its parity, one of $(p, q)$ must be odd and the other must be even. WLOG, assume that $p=2$ and that $q$ is odd. We will multiply by $2$ at the end to take symmetry into account. We have $2^{q}+q^{2}=r$. Consider this equation in mod $3$. The expression $2^{q}$ is $1$ (mod $3$) when $q$ is even and $2$ (mod $3$) when $q$ is odd. However, since $q$ is odd, $2^{q}\equiv2$ (mod $3$). By Fermat's Little Theorem, $q^{2}\equiv1$ (mod $3$) if $q$ is not a multiple of $3$. However, if $q$ is a multiple of $3$, $q^{2}\equiv0$ (mod $3$). If $q$ is not a multiple of $3$, $2+1\equiv0\equiv r$ (mod $3$), which implies $3$ divides $r$. Since $r$ is prime, $r=3$, but $r\ge8$, which is a contradiction. However, if $q$ is a multiple of $3$, we must have $q=3$. This gives $r=17$, which is prime. We have $1$ ordered triple for $p=2$ and multiply by $2$ for symmetry, there are $2$ ordered triples.

\item \textbf{Answer (C):} We wish to determine which of the intervals in the answer choices contains the smallest positive integer $N$ such that $10^{N}>4^{1000}=2^{2000}$. Since $N$ is minimized, we can approximate $10^{N}\approx2^{2000}$. We know that $2^{3}<10<2^{4}$. Raising the inequality to the power of $N$, we get: $2^{3N}<10^{N}<2^{4N}\Rightarrow2^{3N}<2^{2000}<2^{4N}\Rightarrow3N<2000<4N$. $3N<2000\Rightarrow N<667$. $2000<4N\Rightarrow500<N$. Therefore, $500<N<667$. The interval that best matches our bounds on $N$ is $[550,650]$. It turns out that the exact value of $\min(N)$ is $603$.

\item \textbf{Answer (D):} Since $p(a)=2$, $p(p(a))=17\Rightarrow p(2)=17$. Since $p(p(a))=17$, $p(p(p(a)))=167\Rightarrow p(17)=167$. Let $p(x)=mx+n$ for real constants $m$ and $n$. Plugging in $x=2$ gives us $17=2m+n$. Plugging in $x=17$ gives us $167=17m+n$. Solving these equations, $(m,n)=(10,-3)$. Since $p(a)=2$ we have $2=10a-3\Rightarrow a=\frac{1}{2}$.

\item \textbf{Answer (D):} Let $R$ denote a rightwards move and $U$ denote an upwards move. First, arrange all $9$ of the $R$s in a line: $RRRRRRRRR$. We have $10$ spaces to insert a $U$, namely at the beginning of the string, at the end of the string, or in between two $R$s. We can only insert up to $1$ $U$ in each of these $10$ spaces, since each $U$ must be followed by an $R$. In addition, a $U$ as the last character is allowed. There are $\binom{10}{5}=252$ ways to choose the spaces for the $U$ and hence, the number of possible sequences.

\item \textbf{Answer (C):} Let $E^{\prime}$ be the rotation of $E$ $180^{\circ}$ around $D$. Then, $BECE^{\prime}$ is a parallelogram because $BC$ and $EE^{\prime}$ bisect each other. It then follows that $\angle BEC=\angle BE^{\prime}C$. Then, $\angle BEC+\angle BAC=\angle BE^{\prime}C+\angle BAC=180^{\circ}$. It then follows that $BE^{\prime}CA$ is cyclic. By Power of a Point with respect to $D$, $DE^{\prime}\cdot DA=BD\cdot CD$. Because $DE^{\prime}=DE$, we have $DE\cdot DA=BD\cdot CD$. $D$ is the midpoint of $BC$, so we have $BD=CD=\frac{\sqrt{7}}{2}$. Finally, $DE\cdot DA=\frac{7}{4}$. The requested sum is $7+4=11$.

\item \textbf{Answer (D):} Let $k$ be the number of seconds before the expected sum of the outputted numbers is at least $280$. Clearly, one output has no effect on the others, so the expected sum of the outputted numbers is the sum of the expected numbers from each second. At the $n$th second, the expected output is the average of the numbers $\{1,2,3,..., n\}$ or $\frac{n+1}{2}$. Summing this expression from $1$ to $k$, we have: $\frac{1+1}{2}+\frac{2+1}{2}+\frac{3+1}{2}+\cdot\cdot\cdot+\frac{k+1}{2}=\frac{k(k+3)}{4}\ge280\Rightarrow k^{2}+3k-1120\ge0$. By testing values of $k$ or by factoring $k^{2}+3k-1120=(k-32)(k+35)\ge0$, we find that $k=32$ is the smallest value of $k$ that satisfies the inequality.

\item \textbf{Answer (C):} Expanding out the expression, we get $2(ab+ac+ad+bc+bd+cd)$. For this expression to be divisible by $4$, $ab+ac+ad+bc+bd+cd$ must be even. Clearly, $a$, $b$, $c$, and $d$ each have a $\frac{1}{2}$ chance of being odd and a $\frac{1}{2}$ chance of being even. Since $a$, $b$, $c$, and $d$ are all interchangeable in the expression, we do casework on the number of elements in the set $\{a, b, c, d\}$ that are even. There are $2^{4}=16$ possible outcomes. If none of the elements are even, each of the addends are odd, which results in an even expression. There is $1$ way to select $0$ elements. If only $1$ of the elements is even, then $3$ of the addends will be even, while the other $3$ will be odd. This results in an odd expression. If $2$ of the elements are even, then every addend except one of them will be even. The odd addend is the one that is the product of the two odd elements. This results in an odd expression. If $3$ or $4$ of the elements are even, then each addend will be even, which results in an even expression. There are $4$ ways to select $3$ elements and $1$ way to select $4$ elements. Therefore, our probability is $\frac{1+4+1}{16}=\frac{3}{8}$.

\item \textbf{Answer (C):} Let $l_{1}$ be the line $y=2x$ for $x\ge0$ and $l_{2}$ be the line $y=-2x$ for $x\le0$. Clearly, $A$ and $C$ lie on $l_{1}$, while $B$ lies on $l_{2}$. In addition, let $l_{3}$ be the line passing through $M$ that is perpendicular to $l_{1}$. It follows that $B$ is the intersection of $l_{2}$ and $l_{3}$. Since lines $l_{1}$ and $l_{3}$ are perpendicular, the slope of $l_{3}$ is the negative reciprocal of the slope of $l_{1}$. Therefore, the slope of $l_{3}$ is $-\frac{1}{2}$. Since, $l_{3}$ passes through $(3,6)$, the equation of $l_{3}$ is given by $y=-\frac{1}{2}x+\frac{15}{2}$. Since $l_{2}$ has the equation $y=-2x$, $B=(-5,10)$. Because $\triangle ABC$ is equilateral, $AC=BM\cdot\frac{2}{\sqrt{3}}$. Therefore, $[ABC]=\frac{BM^{2}}{\sqrt{3}}$. By the distance formula, $BM^{2}=(3-(-5))^{2}+(6-10)^{2}=80$. Therefore, $[ABC]=\frac{80}{\sqrt{3}}$ or $[ABC]=\frac{80\sqrt{3}}{3}$.

\item \textbf{Answer (C):} Let the path have $a$ southeast arrows, $b$ horizontal arrows, and $c$ southwest arrows. Then, $a+b+c=8$. Since James must travel down $5$ rows, we must have $a+c=5$ as only the southeast and southwest arrows cause James to travel down a row. Lastly, to ensure that James ends up on $AB$ among all the lines parallel to $AB$, we must have $b=c$. Solving these equations gives $(a,b,c)=(2,3,3)$. Let the point at the bottom left corner of the grid be $C$. We must make sure that our path does not go out of the bounds of $\triangle ABC$. Clearly, the directions James can travel in prevent us from ever crossing $AC$. For James to cross $BC$, he must travel down more than $5$ rows. Our composition of directional arrows ensures that he will travel down exactly $5$ rows. Therefore, James can't cross $BC$. However, if at any point in the path, James has traveled more times horizontally than he has in the southwest direction, he will cross $AB$. Consider a string of characters consisting of $3$ $b$s and $3$ $c$s. We will count how many strings there are such that at every point in the string, there are always at least as many $c$s as $b$s. There are $5$ such strings, namely $cbcbcb$, $cbccbb$, $ccbcbb$, $ccbbcb$, and $cccbbb$. The $2$ $a$s can be inserted anywhere in these $5$ strings to construct a possible path. Suppose we have $8$ blank characters in a row. We have $\binom{8}{2}=28$ ways to choose the locations of the $2$ $a$s. Then, we choose any of the $5$ strings of $b$s and $c$s and insert the characters in order in the remaining $6$ spaces. This gives $28\cdot5=140$ paths total.

\item \textbf{Answer (D):} Clearly, $\triangle ABC$ is obtuse in $\angle ABC$ since $AB^{2}+BC^{2}=121+169<400=AC^{2}$. Since $\angle ADC=90^{\circ}$ by the Pythagorean Theorem $(AB+BD)^{2}+CD^{2}=AC^{2}\Rightarrow(11+BD)^{2}+CD^{2}=400$ and $BD^{2}+CD^{2}=BC^{2}=169$. Expanding the first equation, $121+22BD+BD^{2}+CD^{2}=400$. Plugging in the second equation, $121+22BD+169=400\Rightarrow22BD=110\Rightarrow BD=5$. Then, $\triangle ABC \sim \triangle EBD$ because $\angle CBA=\angle DBE$ and $\angle EDB=\angle BCA$ by cyclic quadrilateral $ACDE$. Therefore, $DE=AC\cdot\frac{BD}{BC}=20\cdot\frac{5}{13}=\frac{100}{13}$.

\item \textbf{Answer (B):} We claim that for positive integers $a$, $b$, $x$, and $y$, $a*b=x*y$ if and only if $a+b=x+y$. Our base case for $a+b=x+y=2$ holds, as $(a,b)=(1,1)$ is the only ordered pair of positive integers such that $a+b=2$. Assume that our inductive hypothesis holds true for $a,b,x,$ and $y$. Then, $a*b=x*y$ implies that $(a+1)*b=(x+1)*y$. By the second property of the operation, $(a+1)*b=(x+1)*y=a*b+a+b=x*y+x+y$. By our inductive hypothesis, $a*b+a+b=x*y+x+y$ implies that $a+b=x+y$. This completes the induction for our proposition. In addition, because our steps are reversible, the converse of our proposition holds true as well. Therefore, $a*b=x*y$ if and only if $a+b=x+y$. For a fixed value of $c$, $a*b=c$ must be satisfied by exactly $25$ ordered pairs of positive integers $(a, b)$. Therefore, we must find a positive integer $m$ such that $a+b=m$ has exactly $25$ ordered pairs of positive integer solution. Clearly, $m=26$ and our ordered pairs are $(25,1), (24, 2), (23, 3), ...., (2, 24), (1, 25)$. Therefore, $c=25*1$. Let $f(n)=n*1$ for all $n\ge1$. We are given that $f(1)=2$, and we seek $f(25)$. By the second property of the function, $(n+1)*1=n*1+n+1$ or $f(n+1)=f(n)+(n+1)$. Therefore, $f(25)=2+(2+3+4+...+24+25)=326$.

\item \textbf{Answer (D):} Let $O_{1}$, $O_{2}$, $O_{3}$ be the centers of $\Gamma_{1}$, $\Gamma_{2}$, $\Gamma_{3}$, respectively. By the tangency conditions, $O_{1}O_{2}=2+3=5$, $O_{1}O_{3}=6-2=4$, and $O_{2}O_{3}=6-3=3$, that is to say, $O_{3}$ lies on $\Gamma_{2}$ and since $3^{2}+4^{2}=5^{2}$, $\triangle O_{1}O_{2}O_{3}$ is right in $O_{3}$. Hence, since $O_{3}B=O_{3}C=6$, $\triangle BO_{3}C$ is a right isosceles triangle. Let $O$ be the circumcenter of $\triangle ABC$; by the definition of circumcenter, $OA=OB$, so $\triangle AOB$ is isosceles too. Since $\triangle ABO_{1}$ and $\triangle ACO_{2}$ are isosceles, $\angle BAO_{1}+\angle CAO_{2}=\angle ABO_{1}+\angle ACO_{2}=\frac{1}{2}\angle AO_{1}O_{3}+\frac{1}{2}\angle AO_{2}O_{3}=45^{\circ}$. Hence, $\angle BAC=135^{\circ}$. Finally, by the fact that the angle at the center is twice the angle at the circumference, considering arc $BC$ not containing $A$, the concave $\angle BOC=2\angle BAC=270^{\circ}$, hence $\angle BOC=90^{\circ}$ and $BO_{3}CO$ is a square, that is to say, the circumradius of $\triangle ABC$ is $R=6$.

\item \textbf{Answer (A):} Since $x=\lfloor x\rfloor+\{x\}$ for all real $x$, rewrite the definition of $f(x)$ as $\lfloor x\rfloor(\lfloor x\rfloor+2\{x\})$. Since $0\le\{x\}<1$, we have: $\lfloor x\rfloor^{2}\le f(x)=\lfloor x\rfloor(\lfloor x\rfloor+2\{x\})<\lfloor x\rfloor^{2}+2\lfloor x\rfloor<(\lfloor x\rfloor+1)^{2}$ and $\lfloor x\rfloor^{2}\le f(x)<(\lfloor x\rfloor+1)^{2}$. Since $18^{2}<334<343<19^{2}$, we have $\lfloor\alpha\rfloor=\lfloor\beta\rfloor=18$. In addition, because $\lfloor\alpha\rfloor=\lfloor\beta\rfloor$, $\alpha-\beta=\{\alpha\}-\{\beta\}$. We know that $9=f(\alpha)-f(\beta)=18(18+2\{\alpha\})-18(18+2\{\beta\})$. Then, $9=18(18+2\{\alpha\})-18(18+2\{\beta\})\Rightarrow\frac{1}{4}=\{\alpha\}-\{\beta\}$. Hence, our answer is $\frac{1}{4}$.
\end{enumerate}


\section*{Soal}

\begin{enumerate}

% Answer: D
\item Berapakah nilai dari $201\times9+20\times19-2\times(0+1+9)$?
\begin{enumerate}
    \item $1409$
    \item $1449$
    \item $2019$
    \item $2169$
    \item $2209$
\end{enumerate}

% Answer: E
\item Manakah dari bilangan prima berikut yang merupakan pembagi dari $1^{6}+2^{5}+3^{4}+4^{3}+5^{2}+6^{1}$?
\begin{enumerate}
    \item $2$
    \item $3$
    \item $5$
    \item $7$
    \item $11$
\end{enumerate}

% Answer: A
\item Misalkan $a$, $b$, $c$ dan $d$ adalah bilangan bulat berbeda antara $1$ dan $50$ inklusif. Berapa banyak kemungkinan nilai dari $a+b+c+d$ yang ada?
\begin{enumerate}
    \item $185$
    \item $191$
    \item $196$
    \item $197$
    \item $200$
\end{enumerate}

% Answer: D
\item Misalkan $N$ adalah bilangan bulat positif unik sehingga $N$\% dari $880$ adalah kuadrat sempurna yang lebih kecil dari $880$. Berapakah jumlah digit-digit dari $N$?
\begin{enumerate}
    \item $7$
    \item $8$
    \item $9$
    \item $10$
    \item $11$
\end{enumerate}

% Answer: B
\item Rasio luas segi delapan terhadap kelilingnya adalah $5/4$, ketika perhitungan dilakukan dalam satuan inci (inch). Ketika perhitungan dilakukan dalam kaki (feet), berapakah rasio luas segi delapan yang sama terhadap kelilingnya? (\textit{Catatan: 1 kaki = 12 inci})
\begin{enumerate}
    \item $\frac{5}{576}$
    \item $\frac{5}{48}$
    \item $\frac{5}{4}$
    \item $15$
    \item $180$
\end{enumerate}

% Answer: B
\item Dua bilangan berbeda dipilih secara acak dari himpunan $S=\{1,2,3,4,5,6,7,8,9,10,11,...,2019\}$. Peluang bahwa jumlah keduanya adalah $2020$ akan menjadi lebih besar jika bilangan $n$ terlebih dahulu dihapus dari himpunan $S$. Berapakah jumlah digit-digit dari $n$?
\begin{enumerate}
    \item $1$
    \item $2$
    \item $3$
    \item $10$
    \item $12$
\end{enumerate}

% Answer: D
\item Sepuluh anggota geng (gangster) berdiri di atas permukaan datar dengan jarak antar mereka semuanya berbeda. Pada pukul dua belas, ketika bel gereja mulai berdentang, masing-masing dari mereka menembak mati satu orang di antara $9$ anggota geng lainnya yang jaraknya paling jauh. Paling banyak berapa anggota geng yang akan tetap hidup?
\begin{enumerate}
    \item $1$
    \item $2$
    \item $7$
    \item $8$
    \item $9$
\end{enumerate}

% Answer: D
\item Persegi panjang $ABCD$ digambar dengan $AB=8$ dan $BC=6$. Lingkaran dalam (incircle) dari $\triangle ACD$ memiliki pusat $I$. Berapakah $BI^{2}$?
\begin{enumerate}
    \item $32$
    \item $36$
    \item $48$
    \item $52$
    \item $64$
\end{enumerate}

% Answer: D
\item Kita mendefinisikan bilangan real $x$ sebagai bilangan semi-bulat jika $\sqrt{2}\cdot x$ adalah bilangan bulat. Berapa banyak bilangan real $0\le x\le100$ yang merupakan bilangan semi-bulat?
\begin{enumerate}
    \item $71$
    \item $140$
    \item $141$
    \item $142$
    \item $211$
\end{enumerate}

% Answer: D
\item Berapakah jumlah semua bilangan bulat positif $n$ sehingga $3^{20}<n^{n}<8^{20}$?
\begin{enumerate}
    \item $60$
    \item $65$
    \item $66$
    \item $75$
    \item $84$
\end{enumerate}

% Answer: C
\item Untuk bilangan real $m$ dan $n$, kita mendefinisikan operator $\circ$ sebagai $m\circ n = m^{-1}-(mn)^{-1}+n^{-1}$. Berapakah nilai dari $1\circ(2\circ(\cdot\cdot\cdot(2018\circ2019)))$?
\begin{enumerate}
    \item $\frac{1}{2019}$
    \item $\frac{2018}{2019}$
    \item $1$
    \item $\frac{2019}{2018}$
    \item $2019$
\end{enumerate}

% Answer: B
\item Segi enam beraturan $ABCDEF$ memiliki panjang sisi $1$. Titik $P$ dan $Q$ diletakkan di bagian dalam segi enam sehingga $ABPQ$ adalah sebuah persegi. Jika $M$ dan $N$ masing-masing adalah titik tengah dari sisi $AB$ dan $DE$, berapakah luas daerah yang diarsir?
\begin{enumerate}
    \item $\frac{2}{3}$
    \item $\frac{\sqrt{3}}{2}$
    \item $\frac{8}{9}$
    \item $1$
    \item $\frac{2\sqrt{3}}{3}$
\end{enumerate}

% Answer: C
\item Berapa banyak tripel berurutan bilangan prima $(p,q,r)$ yang memenuhi $p^{q}+q^{p}=r$?
\begin{enumerate}
    \item $0$
    \item $1$
    \item $2$
    \item $3$
    \item $6$
\end{enumerate}

% Answer: C
\item Misalkan $\min(N)$ adalah bilangan bulat positif terkecil sehingga $\frac{0.25^{1000}}{0.1^{N}}>1$. Di antara interval berikut, di manakah $\min(N)$ berada?
\begin{enumerate}
    \item $[250, 350]$
    \item $[400, 500]$
    \item $[550, 650]$
    \item $[700, 800]$
    \item $[850, 950]$
\end{enumerate}

% Answer: D
\item Tinjau polinomial $p(x)$ berderajat $1$ sehingga untuk suatu bilangan real $a$, $p(a)=2$, $p(p(a))=17$ dan $p(p(p(a)))=167$. Berapakah nilai dari $a$?
\begin{enumerate}
    \item $-1$
    \item $-\frac{1}{2}$
    \item $0$
    \item $\frac{1}{2}$
    \item $1$
\end{enumerate}

% Answer: D
\item Geoff si katak berdiri di titik asal pada bidang koordinat. Untuk setiap gerakan, Geoff hanya bisa bergerak $1$ satuan ke kanan atau $1$ satuan ke atas; selain itu, setiap gerakan ke atas harus segera diikuti dengan gerakan ke kanan (kecuali untuk gerakan terakhir). Berapakah jumlah urutan gerakan berbeda yang berakhir di titik $(9,5)$?
\begin{enumerate}
    \item $56$
    \item $126$
    \item $210$
    \item $252$
    \item $346$
\end{enumerate}

% Answer: C
\item Pada segitiga $\triangle ABC$, misalkan $D$ adalah titik tengah $BC$ dan $E$ adalah sebuah titik pada $AD$ sehingga $\angle BEC+\angle BAC=180^{\circ}$. Jika $BC=\sqrt{7}$, nilai dari $DE \cdot DA$ dapat ditulis sebagai $\frac{m}{n}$ untuk suatu bilangan bulat yang saling prima $m$, $n$. Berapakah $m+n$?
\begin{enumerate}
    \item $7$
    \item $9$
    \item $11$
    \item $51$
    \item $53$
\end{enumerate}

% Answer: D
\item Sebuah pembuat bilangan acak menghasilkan sebuah bilangan acak setiap detik. Pada detik ke-$n$, pembuat bilangan tersebut secara acak mengeluarkan sebuah bilangan bulat dari himpunan $\{1,2,..., n\}$. Berapakah ekspektasi jumlah detik yang dibutuhkan sebelum jumlah dari bilangan-bilangan yang dikeluarkan setidaknya mencapai $280$?
\begin{enumerate}
    \item $23$
    \item $24$
    \item $30$
    \item $32$
    \item $35$
\end{enumerate}

% Answer: C
\item Misalkan $a$, $b$, $c$, $d$ adalah $4$ bilangan bulat yang tidak harus berbeda, yang dipilih dari himpunan $\{0,1,..., 2019\}$. Berapakah peluang bahwa $(a+b)(c+d)+(a+c)(b+d)+(a+d)(b+c)$ habis dibagi $4$?
\begin{enumerate}
    \item $\frac{1}{4}$
    \item $\frac{3}{16}$
    \item $\frac{3}{8}$
    \item $\frac{7}{16}$
    \item $\frac{1}{2}$
\end{enumerate}

% Answer: C
\item Misalkan $\triangle ABC$ adalah segitiga sama sisi dengan ketiga titik sudutnya berada pada grafik $y=|2x|$. Jika titik tengah $AC$ adalah $(3,6)$, maka berapakah luas $\triangle ABC$?
\begin{enumerate}
    \item $20\sqrt{3}$
    \item $25\sqrt{3}$
    \item $\frac{80\sqrt{3}}{3}$
    \item $\frac{85\sqrt{3}}{3}$
    \item $45\sqrt{3}$
\end{enumerate}

% Answer: C
\item James harus melakukan perjalanan dari titik $A$ ke titik $B$ di sepanjang kisi (grid) yang ditunjukkan di bawah ini. Dia hanya bisa bergerak mengikuti tanda panah dan dia harus mengambil tepat $8$ langkah. Berapa banyak kemungkinan jalur yang ada?
\begin{enumerate}
    \item $70$
    \item $72$
    \item $140$
    \item $360$
    \item $560$
\end{enumerate}

% Answer: D
\item Segitiga $\triangle ABC$ memiliki $AB=11$, $BC=13$ dan $AC=20$. Sebuah lingkaran digambar dengan diameter $AC$. Garis $AB$ memotong lingkaran di $D\ne A$ dan garis $BC$ memotong lingkaran di $E\ne B$. Berapakah panjang $DE$?
\begin{enumerate}
    \item $\frac{55}{13}$
    \item $5$
    \item $\frac{66}{13}$
    \item $\frac{100}{13}$
    \item $\frac{120}{13}$
\end{enumerate}

% Answer: B
\item Untuk bilangan bulat positif $a$ dan $b$, misalkan $*$ adalah operator dengan sifat-sifat berikut: $a*b=b*a$, $(a+1)*b=a*b+a+b$, $1*1=2$. Misalkan $c$ adalah bilangan bulat positif unik sehingga terdapat tepat $25$ pasangan berurutan bilangan bulat positif $(a, b)$ yang memenuhi $a*b=c$. Berapakah nilai dari $c$?
\begin{enumerate}
    \item $301$
    \item $326$
    \item $351$
    \item $376$
    \item $401$
\end{enumerate}

% Answer: D
\item Misalkan $\Gamma_{1}$, $\Gamma_{2}$, $\Gamma_{3}$ adalah $3$ lingkaran dengan jari-jari berturut-turut $2$, $3$, $6$, sehingga $\Gamma_{1}$ dan $\Gamma_{2}$ bersinggungan luar di $A$ dan $\Gamma_{3}$ menyinggung $\Gamma_{1}$ dan $\Gamma_{2}$ dari dalam berturut-turut di $B$, $C$. Berapakah jari-jari lingkaran luar dari $\triangle ABC$?
\begin{enumerate}
    \item $\frac{6}{\sqrt{11}}$
    \item $5$
    \item $\sqrt{2}+\sqrt{3}+\sqrt{6}$
    \item $6$
    \item $4\sqrt{3}$
\end{enumerate}

% Answer: A
\item Misalkan $f(x)=\lfloor x\rfloor(x+\{x\})$ adalah fungsi yang didefinisikan untuk semua bilangan real positif $x$, di mana $\lfloor x\rfloor$ menyatakan bilangan bulat terbesar yang kurang dari atau sama dengan $x$ dan $\{x\}=x-\lfloor x\rfloor$ menyatakan bagian pecahan dari $x$. Diketahui bahwa terdapat bilangan real unik $\alpha$ dan $\beta$ sehingga $f(\alpha)=343$ dan $f(\beta)=334$. Berapakah $\alpha-\beta$?
\begin{enumerate}
    \item $\frac{1}{4}$
    \item $\frac{1}{3}$
    \item $\frac{3}{8}$
    \item $\frac{5}{12}$
    \item $\frac{4}{9}$
\end{enumerate}

\end{enumerate}

\section*{Solusi}

\begin{enumerate}
\item \textbf{Jawaban (D):} Kita hitung: $201\times9+20\times19-2\times(0+1+9)=1809+380-20=2169$.

\item \textbf{Jawaban (E):} Kita hitung: $1^{6}+2^{5}+3^{4}+4^{3}+5^{2}+6^{1}=1+32+81+64+25+6=209$. Karena $209=11\cdot19$, maka $11$ adalah faktor prima dari ekspresi tersebut.

\item \textbf{Jawaban (A):} Nilai terkecil yang mungkin dari $a+b+c+d$ adalah $1+2+3+4=10$, sedangkan nilai terbesar yang mungkin adalah $47+48+49+50=194$. Tinjau ekspresi $1+2+3+4$. Kita dapat terus menambahkan $4$ dengan $1$ hingga mencapai $1+2+3+50$. Selanjutnya, kita dapat terus menambahkan $3$ dengan $1$ hingga mencapai $1+2+49+50$. Kemudian, kita dapat terus menambahkan $2$ dengan $1$ hingga mencapai $1+48+49+50$. Terakhir, kita dapat terus menambahkan $1$ dengan $1$ hingga mencapai $47+48+49+50$. Oleh karena itu, semua nilai bilangan bulat antara $10$ dan $194$ inklusif dapat diperoleh, sehingga terdapat $194-10+1=185$ kemungkinan nilai.

\item \textbf{Jawaban (D):} Bilangan yang merupakan $N$\% dari $880$ adalah $\frac{N}{100}\cdot880\Rightarrow\frac{44N}{5}$. Karena ekspresi ini harus berupa bilangan bulat, kita misalkan $N=5k$ untuk suatu bilangan bulat positif $k$. Kemudian, ekspresi kita menjadi $44k=2^{2}\cdot11\cdot k$. Agar ini menjadi kuadrat sempurna, eksponen dari setiap bilangan prima dalam faktorisasi prima harus genap. Oleh karena itu, $k=11$ dan $N=55$. Kita dapat memeriksa bahwa $55$\% dari $880$ adalah $484=22^{2}$. Jumlah digit yang diminta adalah $5+5=10$. Untuk menunjukkan bahwa $N$ unik, kita perhatikan bahwa agar bilangan yang dihasilkan lebih kecil dari $880$, kita harus memiliki $N<100$ atau $k<20$. Karena bilangan yang dihasilkan adalah kuadrat sempurna jika $2^{2}\cdot11\cdot k$ adalah kuadrat sempurna, maka $k$ harus habis dibagi $11$. Tetapi $11$ adalah satu-satunya bilangan bulat positif kurang dari $20$ yang habis dibagi $11$, sehingga $N$ tidak bisa bernilai lain yang kurang dari $100$. Jadi, $N$ adalah unik.

\item \textbf{Jawaban (B):} Rasio luas segi delapan terhadap kelilingnya dalam satuan inci adalah $\frac{5~in.^{2}}{4~in.}=\frac{5~in.}{4}$. Kita perhatikan bahwa $1~ft=12~in.$. Mengubah satuan pengukuran ke kaki (feet): $\frac{5~in.}{4}\cdot\frac{1~ft}{12~in.}=\frac{5~ft}{48}$. Oleh karena itu, rasio yang diminta adalah $\frac{5}{48}$.

\item \textbf{Jawaban (B):} Pertama, pilih satu dari $2019$ bilangan secara acak dan kemudian pilih satu dari $2018$ bilangan yang tersisa secara acak. Perhatikan bahwa untuk setiap bilangan bulat positif $k$ dari $1$ hingga $2019$ kecuali $1010$, bilangan $2020-k$ dapat dipilih setelah memilih $k$. Oleh karena itu, jika kita memilih $k\ne1010$ sebagai bilangan pertama, ada peluang $\frac{1}{2018}$ bahwa kita akan memilih $2020-k$ sebagai bilangan kedua. Namun, jika kita memilih $k=1010$ sebagai bilangan pertama, tidak ada peluang kita akan memilih $1010$ lagi, karena kedua bilangan tersebut harus berbeda. Oleh karena itu, $n=1010$ dan jumlah digit yang diminta adalah $1+0+1+0=2$.

\item \textbf{Jawaban (D):} Secara objektif, untuk memaksimalkan jumlah anggota geng yang hidup, satu anggota geng harus ditembak oleh sebanyak mungkin anggota geng untuk mencegah anggota geng lainnya mati. Jelas, ke-$10$ dari mereka tidak bisa tetap hidup. Jika $9$ anggota geng tetap hidup, maka satu anggota geng harus ditembak oleh setiap anggota geng, termasuk dirinya sendiri. Namun, seorang anggota geng jelas tidak bisa menembak dirinya sendiri karena dia bukan anggota geng yang jaraknya paling jauh dari dirinya sendiri. Akan tetapi, $8$ anggota geng bisa tetap hidup. Ini dapat dilakukan dengan menempatkan $9$ anggota geng saling berdekatan dan menempatkan anggota geng ke-$10$ cukup jauh dari anggota geng lainnya. Kemudian, anggota geng ke-$10$ akan menembak salah satu dari $9$ anggota geng pertama, sementara kesembilan anggota geng pertama semuanya akan menembak anggota geng ke-$10$. Oleh karena itu, paling banyak $8$ anggota geng yang bisa tetap hidup.

\item \textbf{Jawaban (D):} Jelas, $AC=10$ berdasarkan Teorema Pythagoras pada $\triangle ACD$. Tinjau lingkaran dalam dari $\triangle ACD$. Jari-jarinya sama dengan luas $\triangle ACD$ dibagi dengan semi-kelilingnya (setengah keliling). Luas $\triangle ACD$ adalah $\frac{1}{2}\cdot6\cdot8=24$. Karena semi-kelilingnya adalah $12$, jari-jari lingkaran dalamnya adalah $2$. Oleh karena itu, jarak dari $I$ ke $DC$ adalah $2$ dan jarak dari $I$ ke $AB$ adalah $6-2=4$. Jarak dari $I$ ke $AD$ adalah $2$, sehingga jarak dari $I$ ke $BC$ adalah $8-2=6$. Oleh karena itu, berdasarkan Teorema Pythagoras, $BI^{2}=4^{2}+6^{2}=52$.

\item \textbf{Jawaban (D):} Agar $\sqrt{2}\cdot x$ menjadi bilangan bulat, $x$ harus berbentuk $\frac{n}{\sqrt{2}}$ untuk suatu bilangan bulat $n$. Maka, $0\le x\le100\Rightarrow0\le\frac{n}{\sqrt{2}}\le100\Rightarrow0\le n\le100\sqrt{2}$. Nilai $\sqrt{2}$ sering diperkirakan sebagai $1.414$, sehingga $141<100\sqrt{2}<142\Rightarrow0\le n\le141$. Ada $142$ kemungkinan nilai $n$, yang memberi kita $142$ kemungkinan nilai $x$.

\item \textbf{Jawaban (D):} Perhatikan bahwa $3^{20}=9^{10}$ dan $8^{20}=16^{15}$. Ketidaksamaan kita menjadi $9^{10}<n^{n}<16^{15}$. Jelas, nilai-nilai $n$ yang memenuhi ketidaksamaan ini adalah $10, 11, 12, 13, 14$, dan $15$, yang jumlahnya adalah $75$.

\item \textbf{Jawaban (C):} Perhatikan bahwa untuk semua bilangan real $k$, $1\circ k=1-\frac{1}{k}+\frac{1}{k}=1$. Oleh karena itu, $1\circ(2\circ(\cdot\cdot\cdot(2018\circ2019)))=1$.

\item \textbf{Jawaban (B):} Luas daerah yang diarsir diberikan oleh luas layang-layang $NCMF$ dikurangi luas layang-layang $NPMQ$. Karena luas layang-layang adalah setengah dari hasil kali diagonal-diagonalnya, luas yang kita inginkan adalah $\frac{1}{2}(NM\cdot CF-NM\cdot PQ)$. Jelas, $NM=\sqrt{3}$ dan $CF=2$. Selain itu, $PQ=1$ karena $ABPQ$ adalah sebuah persegi. Oleh karena itu, luas yang diinginkan adalah $\frac{1}{2}(\sqrt{3}\cdot2-\sqrt{3}\cdot1)=\frac{\sqrt{3}}{2}$.

\item \textbf{Jawaban (C):} Perhatikan bahwa $r=p^{q}+q^{p}\ge2^{2}+2^{2}=8$. Oleh karena itu, $r\ne2$, sehingga $r$ harus berupa bilangan prima ganjil. Karena pemangkatan suatu bilangan bulat mempertahankan paritasnya (keganjilan/kegenapannya), salah satu dari $(p, q)$ harus ganjil dan yang lainnya harus genap. Tanpa mengurangi keumuman (WLOG), asumsikan bahwa $p=2$ dan $q$ ganjil. Kita akan mengalikan dengan $2$ di akhir untuk memperhitungkan simetri. Kita memiliki $2^{q}+q^{2}=r$. Tinjau persamaan ini dalam modulo $3$. Ekspresi $2^{q}$ kongruen $1$ (mod $3$) ketika $q$ genap dan $2$ (mod $3$) ketika $q$ ganjil. Namun, karena $q$ ganjil, $2^{q}\equiv2$ (mod $3$). Berdasarkan Teorema Kecil Fermat, $q^{2}\equiv1$ (mod $3$) jika $q$ bukan kelipatan $3$. Namun, jika $q$ adalah kelipatan $3$, $q^{2}\equiv0$ (mod $3$). Jika $q$ bukan kelipatan $3$, $2+1\equiv0\equiv r$ (mod $3$), yang mengimplikasikan $3$ habis membagi $r$. Karena $r$ prima, $r=3$, tetapi $r\ge8$, yang merupakan kontradiksi. Namun, jika $q$ adalah kelipatan $3$, kita harus memiliki $q=3$. Ini memberikan $r=17$, yang merupakan bilangan prima. Kita memiliki $1$ tripel berurutan untuk $p=2$ dan mengalikannya dengan $2$ untuk simetri, terdapat $2$ tripel berurutan.

\item \textbf{Jawaban (C):} Kita ingin menentukan interval mana di pilihan jawaban yang memuat bilangan bulat positif terkecil $N$ sehingga $10^{N}>4^{1000}=2^{2000}$. Karena $N$ diminimalkan, kita dapat memperkirakan $10^{N}\approx2^{2000}$. Kita tahu bahwa $2^{3}<10<2^{4}$. Memangkatkan ketidaksamaan ke pangkat $N$, kita dapatkan: $2^{3N}<10^{N}<2^{4N}\Rightarrow2^{3N}<2^{2000}<2^{4N}\Rightarrow3N<2000<4N$. $3N<2000\Rightarrow N<667$. $2000<4N\Rightarrow500<N$. Oleh karena itu, $500<N<667$. Interval yang paling sesuai dengan batas kita untuk $N$ adalah $[550,650]$. Ternyata nilai eksak dari $\min(N)$ adalah $603$.

\item \textbf{Jawaban (D):} Karena $p(a)=2$, $p(p(a))=17\Rightarrow p(2)=17$. Karena $p(p(a))=17$, $p(p(p(a)))=167\Rightarrow p(17)=167$. Misalkan $p(x)=mx+n$ untuk konstanta real $m$ dan $n$. Mensubstitusikan $x=2$ memberi kita $17=2m+n$. Mensubstitusikan $x=17$ memberi kita $167=17m+n$. Menyelesaikan persamaan-persamaan ini, $(m,n)=(10,-3)$. Karena $p(a)=2$ kita peroleh $2=10a-3\Rightarrow a=\frac{1}{2}$.

\item \textbf{Jawaban (D):} Misalkan $R$ menyatakan gerakan ke kanan dan $U$ menyatakan gerakan ke atas. Pertama, susun semua $9$ buah $R$ dalam satu baris: $RRRRRRRRR$. Kita memiliki $10$ ruang untuk menyisipkan $U$, yaitu di awal baris, di akhir baris, atau di antara dua $R$. Kita hanya dapat menyisipkan hingga $1$ buah $U$ di masing-masing dari $10$ ruang ini, karena setiap $U$ harus diikuti oleh $R$. Selain itu, sebuah $U$ sebagai karakter terakhir diperbolehkan. Terdapat $\binom{10}{5}=252$ cara untuk memilih ruang bagi $U$ dan oleh karena itu, terdapat sebanyak itu urutan yang mungkin.

\item \textbf{Jawaban (C):} Misalkan $E^{\prime}$ adalah rotasi $E$ sebesar $180^{\circ}$ dengan pusat $D$. Maka, $BECE^{\prime}$ adalah sebuah jajargenjang karena $BC$ dan $EE^{\prime}$ saling membagi dua sama panjang. Akibatnya $\angle BEC=\angle BE^{\prime}C$. Kemudian, $\angle BEC+\angle BAC=\angle BE^{\prime}C+\angle BAC=180^{\circ}$. Hal ini mengakibatkan segi empat $BE^{\prime}CA$ adalah siklik (segi empat tali busur). Berdasarkan Kuasa Titik terhadap $D$, $DE^{\prime}\cdot DA=BD\cdot CD$. Karena $DE^{\prime}=DE$, kita peroleh $DE\cdot DA=BD\cdot CD$. $D$ adalah titik tengah $BC$, sehingga kita memiliki $BD=CD=\frac{\sqrt{7}}{2}$. Akhirnya, $DE\cdot DA=\frac{7}{4}$. Jumlah yang diminta adalah $7+4=11$.

\item \textbf{Jawaban (D):} Misalkan $k$ adalah jumlah detik sebelum ekspektasi jumlah dari bilangan yang dikeluarkan setidaknya mencapai $280$. Jelas, satu keluaran tidak berpengaruh pada keluaran lainnya, sehingga ekspektasi jumlah dari bilangan yang dikeluarkan adalah jumlah dari nilai ekspektasi setiap detiknya. Pada detik ke-$n$, ekspektasi keluarannya adalah rata-rata dari bilangan-bilangan $\{1,2,3,..., n\}$ yaitu $\frac{n+1}{2}$. Menjumlahkan ekspresi ini dari $1$ hingga $k$, kita peroleh: $\frac{1+1}{2}+\frac{2+1}{2}+\frac{3+1}{2}+\cdot\cdot\cdot+\frac{k+1}{2}=\frac{k(k+3)}{4}\ge280\Rightarrow k^{2}+3k-1120\ge0$. Dengan menguji nilai-nilai $k$ atau dengan memfaktorkan $k^{2}+3k-1120=(k-32)(k+35)\ge0$, kita temukan bahwa $k=32$ adalah nilai terkecil dari $k$ yang memenuhi ketidaksamaan tersebut.

\item \textbf{Jawaban (C):} Dengan menjabarkan ekspresi tersebut, kita peroleh $2(ab+ac+ad+bc+bd+cd)$. Agar ekspresi ini habis dibagi $4$, $ab+ac+ad+bc+bd+cd$ harus genap. Jelas, $a$, $b$, $c$, dan $d$ masing-masing memiliki peluang $\frac{1}{2}$ untuk bernilai ganjil dan peluang $\frac{1}{2}$ untuk genap. Karena $a$, $b$, $c$, dan $d$ semuanya dapat saling dipertukarkan dalam ekspresi tersebut, kita membagi kasus berdasarkan jumlah elemen dalam himpunan $\{a, b, c, d\}$ yang bernilai genap. Terdapat $2^{4}=16$ kemungkinan hasil. Jika tidak ada elemen yang genap, setiap suku penjumlahan adalah ganjil, yang menghasilkan ekspresi genap. Ada $1$ cara untuk memilih $0$ elemen. Jika hanya $1$ elemen yang genap, maka $3$ suku penjumlahan akan genap, sedangkan $3$ suku lainnya akan ganjil. Ini menghasilkan ekspresi ganjil. Jika $2$ elemen genap, maka setiap suku penjumlahan kecuali satu dari mereka akan genap. Suku ganjil tersebut adalah hasil kali dari dua elemen ganjil. Ini menghasilkan ekspresi ganjil. Jika $3$ atau $4$ elemen genap, maka setiap suku penjumlahan akan genap, yang menghasilkan ekspresi genap. Ada $4$ cara untuk memilih $3$ elemen dan $1$ cara untuk memilih $4$ elemen. Oleh karena itu, peluangnya adalah $\frac{1+4+1}{16}=\frac{3}{8}$.

\item \textbf{Jawaban (C):} Misalkan $l_{1}$ adalah garis $y=2x$ untuk $x\ge0$ dan $l_{2}$ adalah garis $y=-2x$ untuk $x\le0$. Jelas, $A$ dan $C$ terletak pada $l_{1}$, sedangkan $B$ terletak pada $l_{2}$. Selain itu, misalkan $l_{3}$ adalah garis yang melalui $M$ dan tegak lurus terhadap $l_{1}$. Mengakibatkan $B$ adalah perpotongan dari $l_{2}$ dan $l_{3}$. Karena garis $l_{1}$ dan $l_{3}$ saling tegak lurus, gradien $l_{3}$ adalah negatif kebalikan dari gradien $l_{1}$. Oleh karena itu, gradien $l_{3}$ adalah $-\frac{1}{2}$. Karena $l_{3}$ melalui $(3,6)$, persamaan garis $l_{3}$ diberikan oleh $y=-\frac{1}{2}x+\frac{15}{2}$. Karena $l_{2}$ memiliki persamaan $y=-2x$, titik $B=(-5,10)$. Karena $\triangle ABC$ adalah segitiga sama sisi, $AC=BM\cdot\frac{2}{\sqrt{3}}$. Oleh karena itu, luas $[ABC]=\frac{BM^{2}}{\sqrt{3}}$. Menggunakan rumus jarak, $BM^{2}=(3-(-5))^{2}+(6-10)^{2}=80$. Oleh karena itu, $[ABC]=\frac{80}{\sqrt{3}}$ atau $[ABC]=\frac{80\sqrt{3}}{3}$.

\item \textbf{Jawaban (C):} Misalkan jalur tersebut memiliki $a$ panah arah tenggara, $b$ panah arah horizontal, dan $c$ panah arah barat daya. Maka, $a+b+c=8$. Karena James harus turun $5$ baris, kita harus memiliki $a+c=5$ karena hanya panah tenggara dan barat daya yang menyebabkan James turun satu baris. Terakhir, untuk memastikan bahwa James berakhir di $AB$ di antara semua garis yang sejajar dengan $AB$, kita harus memiliki $b=c$. Menyelesaikan persamaan-persamaan ini memberikan $(a,b,c)=(2,3,3)$. Misalkan titik di pojok kiri bawah kisi adalah $C$. Kita harus memastikan bahwa jalur kita tidak keluar dari batas $\triangle ABC$. Jelas, arah yang dapat dilalui James mencegah kita untuk menyeberangi $AC$. Agar James menyeberangi $BC$, dia harus turun lebih dari $5$ baris. Komposisi panah arah kita memastikan bahwa ia akan turun tepat $5$ baris. Oleh karena itu, James tidak bisa menyeberangi $BC$. Namun, jika pada titik mana pun di jalurnya, James telah bergerak secara horizontal lebih banyak daripada pergerakannya ke arah barat daya, dia akan menyeberangi $AB$. Tinjau sebuah untaian (string) karakter yang terdiri dari $3$ buah $b$ dan $3$ buah $c$. Kita akan menghitung berapa banyak untaian yang ada sehingga pada setiap titik dalam untaian tersebut, selalu terdapat setidaknya sejumlah $c$ yang sama banyaknya dengan $b$. Terdapat $5$ untaian seperti itu, yaitu $cbcbcb$, $cbccbb$, $ccbcbb$, $ccbbcb$, dan $cccbbb$. Ke-$2$ buah $a$ dapat disisipkan di mana saja dalam $5$ untaian ini untuk membentuk jalur yang mungkin. Misalkan kita memiliki $8$ karakter kosong berturut-turut. Kita memiliki $\binom{8}{2}=28$ cara untuk memilih lokasi ke-$2$ buah $a$ tersebut. Kemudian, kita memilih salah satu dari $5$ untaian $b$ dan $c$ lalu menyisipkan karakter-karakter tersebut secara berurutan di $6$ ruang yang tersisa. Ini memberikan total $28\cdot5=140$ jalur.

\item \textbf{Jawaban (D):} Jelas, $\triangle ABC$ adalah segitiga tumpul di $\angle ABC$ karena $AB^{2}+BC^{2}=121+169<400=AC^{2}$. Karena $\angle ADC=90^{\circ}$ berdasarkan Teorema Pythagoras $(AB+BD)^{2}+CD^{2}=AC^{2}\Rightarrow(11+BD)^{2}+CD^{2}=400$ dan $BD^{2}+CD^{2}=BC^{2}=169$. Menjabarkan persamaan pertama, $121+22BD+BD^{2}+CD^{2}=400$. Mensubstitusi persamaan kedua, $121+22BD+169=400\Rightarrow22BD=110\Rightarrow BD=5$. Kemudian, $\triangle ABC \sim \triangle EBD$ karena $\angle CBA=\angle DBE$ dan $\angle EDB=\angle BCA$ berdasarkan segi empat tali busur $ACDE$. Oleh karena itu, $DE=AC\cdot\frac{BD}{BC}=20\cdot\frac{5}{13}=\frac{100}{13}$.

\item \textbf{Jawaban (B):} Kita asumsikan bahwa untuk bilangan bulat positif $a$, $b$, $x$, dan $y$, berlaku $a*b=x*y$ jika dan hanya jika $a+b=x+y$. Kasus dasar kita untuk $a+b=x+y=2$ berlaku, karena $(a,b)=(1,1)$ adalah satu-satunya pasangan berurutan bilangan bulat positif yang memenuhi $a+b=2$. Asumsikan hipotesis induksi kita berlaku untuk $a,b,x,$ dan $y$. Maka, $a*b=x*y$ mengimplikasikan bahwa $(a+1)*b=(x+1)*y$. Berdasarkan sifat kedua dari operasi tersebut, $(a+1)*b=(x+1)*y=a*b+a+b=x*y+x+y$. Berdasarkan hipotesis induksi kita, $a*b+a+b=x*y+x+y$ mengimplikasikan bahwa $a+b=x+y$. Ini melengkapi induksi untuk proposisi kita. Selain itu, karena langkah-langkah kita dapat dibalik, kebalikan dari proposisi kita juga berlaku. Oleh karena itu, $a*b=x*y$ jika dan hanya jika $a+b=x+y$. Untuk nilai $c$ yang tetap, $a*b=c$ harus dipenuhi oleh tepat $25$ pasangan berurutan bilangan bulat positif $(a, b)$. Oleh karena itu, kita harus menemukan bilangan bulat positif $m$ sehingga $a+b=m$ memiliki tepat $25$ solusi pasangan berurutan bilangan bulat positif. Jelas, $m=26$ dan pasangan berurutan kita adalah $(25,1), (24, 2), (23, 3), ...., (2, 24), (1, 25)$. Oleh karena itu, $c=25*1$. Misalkan $f(n)=n*1$ untuk semua $n\ge1$. Kita diberi tahu bahwa $f(1)=2$, dan kita mencari $f(25)$. Berdasarkan sifat kedua dari fungsi, $(n+1)*1=n*1+n+1$ atau $f(n+1)=f(n)+(n+1)$. Oleh karena itu, $f(25)=2+(2+3+4+...+24+25)=326$.

\item \textbf{Jawaban (D):} Misalkan $O_{1}$, $O_{2}$, $O_{3}$ berturut-turut adalah pusat dari $\Gamma_{1}$, $\Gamma_{2}$, $\Gamma_{3}$. Berdasarkan kondisi persinggungan, $O_{1}O_{2}=2+3=5$, $O_{1}O_{3}=6-2=4$, dan $O_{2}O_{3}=6-3=3$, dengan kata lain, $O_{3}$ terletak pada $\Gamma_{2}$ dan karena $3^{2}+4^{2}=5^{2}$, $\triangle O_{1}O_{2}O_{3}$ siku-siku di $O_{3}$. Sehingga, karena $O_{3}B=O_{3}C=6$, $\triangle BO_{3}C$ adalah segitiga siku-siku sama kaki. Misalkan $O$ adalah pusat lingkaran luar $\triangle ABC$; berdasarkan definisi pusat lingkaran luar, $OA=OB$, sehingga $\triangle AOB$ juga sama kaki. Karena $\triangle ABO_{1}$ dan $\triangle ACO_{2}$ sama kaki, $\angle BAO_{1}+\angle CAO_{2}=\angle ABO_{1}+\angle ACO_{2}=\frac{1}{2}\angle AO_{1}O_{3}+\frac{1}{2}\angle AO_{2}O_{3}=45^{\circ}$. Sehingga, $\angle BAC=135^{\circ}$. Terakhir, berdasarkan fakta bahwa sudut pusat adalah dua kali sudut keliling, dengan meninjau busur $BC$ yang tidak memuat $A$, sudut refleks $\angle BOC=2\angle BAC=270^{\circ}$, sehingga sudut $\angle BOC=90^{\circ}$ dan $BO_{3}CO$ adalah sebuah persegi, dengan kata lain, jari-jari lingkaran luar $\triangle ABC$ adalah $R=6$.

\item \textbf{Jawaban (A):} Karena $x=\lfloor x\rfloor+\{x\}$ untuk semua bilangan real $x$, tulis ulang definisi $f(x)$ sebagai $\lfloor x\rfloor(\lfloor x\rfloor+2\{x\})$. Karena $0\le\{x\}<1$, kita peroleh: $\lfloor x\rfloor^{2}\le f(x)=\lfloor x\rfloor(\lfloor x\rfloor+2\{x\})<\lfloor x\rfloor^{2}+2\lfloor x\rfloor<(\lfloor x\rfloor+1)^{2}$ dan $\lfloor x\rfloor^{2}\le f(x)<(\lfloor x\rfloor+1)^{2}$. Karena $18^{2}<334<343<19^{2}$, kita peroleh $\lfloor\alpha\rfloor=\lfloor\beta\rfloor=18$. Selain itu, karena $\lfloor\alpha\rfloor=\lfloor\beta\rfloor$, berlaku $\alpha-\beta=\{\alpha\}-\{\beta\}$. Kita tahu bahwa $9=f(\alpha)-f(\beta)=18(18+2\{\alpha\})-18(18+2\{\beta\})$. Maka, $9=18(18+2\{\alpha\})-18(18+2\{\beta\})\Rightarrow\frac{1}{4}=\{\alpha\}-\{\beta\}$. Jadi, jawaban kita adalah $\frac{1}{4}$.
\end{enumerate}

\end{document}